\subsection{四項VK圏の内在的区別と次の対象}

$m=3$ の二つの非零CGG量子化について、四項圏そのものの導来非同値が確定した。
外側の補完環の差を使う必要はない。
$HH^1$ は両方とも一次元なので、非零生成元 $\xi$ の作用
$[\xi,-]$ の階数と冪零Jordan分割が、生成元の選択によらない導来不変量になる。
$\alpha=0,s=1$ の代表で $\xi=[\log(1+N)]$ を使うと
\[
\begin{array}{c|c|c}
(j,r)&HH^2\text{ のJordan分割}&HH^3\text{ のJordan分割}\\\hline
(0,3)&(13,10,7,5)&(14,11,8,5,1)\\
(1,1)&(11,9,6,5,3,1)&(13,10,7,5,3,1)
\end{array}
\]
を得る。従って一次変形の固定方向は四次元と六次元である。
可換点も含めれば、この明示CGG族の四項VK圏にはちょうど三つの導来同値類がある。
どちらの非零点でも通常のhullは三十五次元の滑らかな環なので、
hullとHH次元が一致することは圏の同値を意味しない。

再利用できる計算法は、全cocycle空間の基底行列 $K$ を使う式
\[
 \operatorname{rank}(L:HH^2\to HH^2)
 =\operatorname{rank}[d^1\mid LK]-\operatorname{rank}d^1
\]
である。商座標による計算から独立に、第一階数だけで今回の非同値性を確認できる。
導分には $N$ でなく $\log(1+N)$ を使う。
補助次数 $(0,1,2,m+r+1)$ に対する下降作用が、有限複体上の冪零性を保証する。
反対代数への入力反転はこの作用を共役するので、End合成の規約によらない。

次に調べる価値があるのは、全ての $HH^2$ 方向の自己同型作用と、
標準補完方向がその中で占める位置の関係である。
一次の固定部分空間だけでなく、外的導分を持ち上げる条件の全完備局所環も決定した。
全次数の対称性保持族、純粋部分への還元、標準補完接線の二次消失を後述する。
一つの外部導分が与える作用は一般のVK変形へ移ると消滅し得るため、
共鳴点以外の層を分類するには、自己同型群とHochschild括弧の同時変化を追う必要がある。


全相対Hochschild dg Lie代数 $C_E^\bullet(B)[1]$ は、
どちらの非零成分でもhomotopy abelianではない。
$HH^1\times HH^2\to HH^2$ の非零括弧がその証拠であり、
異なる作用階数は二つのdg Lie代数がquasi-isomorphicでないことも示す。
これは三十五次元の滑らかなhullと両立する。
ここで全複体の $C_E^0$ は次数 $-1$ に残す。
$C_E^{\ge1}[1]$ への切捨てでは $H^0$ が $\operatorname{Der}_E(B)$ となり、$HH^1$ との混同が生じる。

\paragraph{元の射空間に見えるJordan鎖.}
$V_2=\langle x,y,z\rangle$ に $N_0=x\partial_y+2y\partial_z$ を作用させると、
通常の次数 $m$ の多項式空間はJordan鎖の直和
\[
 \mathbb C[x,y,z]_m
   =\bigoplus_{k=0}^{\lfloor m/2\rfloor}q^k V_{2(m-2k)}
\]
に分解する。ここで $V_{2r}=\langle N_0^a z^r:0\le a\le2r\rangle$ とする。
その鎖の長さは $2m-4k+1$ である。
$Nw=-q^jz^{m-2j}$ を加えると、$k=j$ の鎖だけが一つ長くなり、
長さ $2(m-2j)+1$ が $2(m-2j)+2$ へ変わる。
$\delta=N(1-N/2+\cdots)$ は $N$ と同じJordan分割を持つ。
従って $m=3$ の次数三射空間での分割は $(8,3)$ と $(7,4)$ となる。

これは普通の代数非同型の短い確認にも使える。
四頂点の到達関係は全順序なので、代数同型を内的同型で補正すれば各冪等元を固定できる。
そのとき $\operatorname{Der}_E(B)=\mathbb C\delta+\operatorname{ad}(E)$ であり、
$\delta$ の像は $c\delta+\operatorname{ad}(h)$（$c\ne0,h\in E$）である。
各射空間上では後項は $h_i-h_j$ というスカラーなので、
像が冪零であることから全差が零となり、像は $c\delta$ に限られる。
次数三射空間のJordan分割の差から代数非同型が従う。
この議論は三角同値には直接適用できないため、導来非同値には本文のHochschild作用を使う。
一般重みの四項cornerで同じ議論を内在化するには、まず $HH^1$ の一次元性が必要である。
